\begin{abstract}
Resolving which cell states are present in a tissue is essential for characterizing cellular heterogeneity in development and disease. Single-cell and spatial transcriptomics profile many donors and sections at once, but individual cells are too sparse to analyze alone, so cells are aggregated into metacells. Current methods form metacells from expression proximity alone: a state rare in every batch but common overall cannot be formed, and a cell's spatial niche is either absorbed into cell-type signal or, if added to the affinity directly, overcorrects and splits one state across neighborhoods. We introduce scProto, a prototype-based deep clustering model that learns metacells across an entire experiment rather than within each batch, using the affinity graph as a training objective rather than clustering it directly: a conditional variational autoencoder learns a batch-corrected representation jointly with soft prototype assignment under a cohesion criterion that keeps each group internally dense via both direct and shared-neighborhood similarity. Assignments come from expression alone, so a niche separates cells only to the extent it changes their transcriptome. Across single-cell and spatial datasets, scProto recovers rare and niche-associated states missed by within-batch clustering, improving batch mixing and rare-population recovery while preserving neighborhood structure.
\end{abstract}
